% Generated draft: requires mathematical and source review.
\documentclass[11pt]{article}
\usepackage[margin=25mm]{geometry}
\usepackage{fontspec}
\setmainfont{Latin Modern Roman}
\usepackage{amsmath,amssymb,mathtools,unicode-math}
\setmathfont{Latin Modern Math}
\usepackage{enumitem,needspace,xurl,hyperref}
\setlength{\parindent}{0pt}
\setlength{\parskip}{5pt}
\setlength{\emergencystretch}{3em}
\newcommand{\N}{\mathbb{N}}
\newcommand{\Z}{\mathbb{Z}}
\newcommand{\Q}{\mathbb{Q}}
\newcommand{\R}{\mathbb{R}}
\newcommand{\C}{\mathbb{C}}
\newcommand{\F}{\mathbb{F}}
\newcommand{\A}{\mathbb{A}}
\newcommand{\PP}{\mathbb{P}}
\newcommand{\E}{\mathbb{E}}
\newcommand{\eps}{\varepsilon}
\newcommand{\dd}{\,\mathrm{d}}
\begin{document}
\section*{UnsolvedMath ID: 30005719}
\textbf{Source problem code:} catalogue entry
\subsection{Definitions and mathematical statement}
For reference, use the following standard minimum convention; the supplied record does not specify its intended domain or sign convention. Let \(d\geq 1\), let \(e_1,\ldots,e_d\) be the standard basis vectors of \(\Z^d\), and let \(f:\Z^d\to\R\cup\{+\infty\}\) have nonempty effective domain \(\operatorname{dom}(f)=\{x:f(x)<+\infty\}\). The function \(f\) is \(M\)-convex if, for every \(x,y\in\operatorname{dom}(f)\) and every \(i\) with \(x_i>y_i\), there exists \(j\) with \(x_j<y_j\) such that
\[
f(x)+f(y)\geq f(x-e_i+e_j)+f(y+e_i-e_j).
\]
A nonempty set \(B\subseteq\Z^d\) is \(M\)-convex when its extended indicator, equal to \(0\) on \(B\) and \(+\infty\) elsewhere, is \(M\)-convex. In this convention, a valuated matroid is an \(M\)-convex function whose effective domain consists of incidence vectors of subsets of a fixed cardinality of \(\{1,\ldots,d\}\).

For a field \(K\), a finite subspace arrangement means a tuple \((V_1,\ldots,V_d)\) of linear subspaces of a finite-dimensional \(K\)-vector space \(W\). As a reference for the requested extension, if \(U\subseteq K^m\) has dimension \(r\), choose an \(r\times m\) matrix \(A\) whose rows form a basis of \(U\). Its ordinary Pl\"ucker coordinates are
\[
p_I=\det(A_I),\qquad I\subseteq\{1,\ldots,m\},\quad |I|=r,
\]
where columns are ordered by their indices. Changing the row basis multiplies all coordinates by the same nonzero scalar. Over a field with a real-valued valuation, valuations of these minors, with the valuation of zero taken to be \(+\infty\), give the usual reference notion of a realizable valuated matroid.

\textbf{Requested definition problem.} Give a natural definition of a \emph{realizable \(M\)-convex function} and of a \emph{Pl\"ucker vector of a subspace arrangement}, extending the corresponding notions for valuated matroids and \(M\)-convex sets. The supplied record does not specify the ground field, valuation, admissible function domains, coordinate index set, equivalence relation on the proposed vectors, or the precise specialization requirements. The reference conventions above therefore do not resolve the intended generalization. This is a request for definitions, with no supplied axioms sufficient to turn it into a uniquely specified existence assertion.

\textbf{Unresolved source qualification.} The original extraction also contains the heading ``Is log-concavity arising through matrix recursion?'' and assertions about undefined objects \(T_n,Q_n,C_n,V_n\): for \(n\leq300\), these are reported to be ultra-log-concave, with Lorentzian homogenizations, while in general they are neither symmetric nor real-rooted. The accompanying literature assessment concerns an all-\(n\) ultra-log-concavity conjecture for a recursive sequence. No recurrence or definitions of these objects are supplied, and their relationship to the realizability question is unspecified. These assertions cannot be incorporated as hypotheses or conclusions of the requested definition problem.
\subsection{Short English statement}
The problem asks for a geometric meaning of realizability for general \(M\)-convex functions and for an analogue of Pl\"ucker coordinates for arrangements of subspaces. The intended compatibility conditions and underlying field are unspecified. The record also contains apparently unrelated material about log-concavity, so the formulation requires human mathematical and source review.
\subsection{Sources}
\begin{itemize}
\item Local UnsolvedMath registry, permanent ID \(30005719\), title \emph{Realizability of M-Convex Functions}, exact source problem code \texttt{catalogue entry}. The supplied record labels the problem open and its statement \texttt{corrected\_verified}, and reports reconstruction and checking against the official report. These are supplied registry assertions, not independent verification here. No proposer is named; the supplied proposed year is 2023.
\item Supplied report citation: \emph{Geometric, Algebraic and Topological Combinatorics} (2024), Oberwolfach Reports, pp.~3249--3318, DOI \texttt{10.4171/owr/2023/58}. \url{https://doi.org/10.4171/owr/2023/58}. This same URL is supplied as the source URL, the background report link, and the literature source labeled ``original Oberwolfach report.'' No exact location within the report is supplied for this question.
\item Reference [1] in the extracted text: L.~Lov\'asz, \emph{Flats in matroids and geometric graphs}, Comb. Surv., Proc. 6th Br. comb. Conf. (1977), pp.~45--86. Bibliographic details are reproduced from the supplied record.
\item Reference [2] in the extracted text: K.~Murota, \emph{Discrete Convex Analysis}, SIAM Monographs on Discrete Mathematics and Applications, Society for Industrial and Applied Mathematics, Philadelphia, PA, 2003.
\item Supplied literature assessment dated exactly \texttt{2026-08-22}: the all-\(n\) ultra-log-concavity assertion for a recursive sequence is described as a conjecture, computationally checked only through \(n\leq300\) in the report, with no general proof identified by that assessment. Its supplied checked source is \url{https://doi.org/10.4171/owr/2023/58}. This is supplied dated curation, not independently verified; the background explicitly requests independent verification of time-sensitive status and cited claims. Its subject does not match the realizability question.
\end{itemize}
\end{document}
